%-------------------------
% Resume in Latex
% Author : Sourabh Bajaj
% License : MIT
%------------------------

\documentclass[letterpaper,11pt]{article}

\usepackage{latexsym}
\usepackage[empty]{fullpage}
\usepackage{titlesec}
\usepackage{marvosym}
\usepackage[usenames,dvipsnames]{color}
\usepackage{enumitem}
\usepackage[pdftex]{hyperref}
\usepackage{fancyhdr}

\pagestyle{fancy}
\fancyhf{}
\fancyfoot{}
\renewcommand{\headrulewidth}{0pt}
\renewcommand{\footrulewidth}{0pt}

% Adjust margins
\addtolength{\oddsidemargin}{-0.375in}
\addtolength{\evensidemargin}{-0.375in}
\addtolength{\textwidth}{1in}
\addtolength{\topmargin}{-.5in}
\addtolength{\textheight}{1.0in}

\urlstyle{same}
\raggedbottom
\raggedright
\setlength{\tabcolsep}{0in}

% Sections formatting
\titleformat{\section}{
  \vspace{-4pt}\scshape\raggedright\large
}{}{0em}{}[\color{black}\titlerule \vspace{-5pt}]

%-------------------------
% Custom commands
\newcommand{\resumeItem}[2]{
  \item\small{
    \textbf{#1}{: #2 \vspace{-3pt}}
  }
}

\newcommand{\resumeSubheading}[4]{
  \vspace{-2pt}\item
    \begin{tabular*}{0.97\textwidth}{l@{\extracolsep{\fill}}r}
      \textbf{#1} & #2 \\
      \textit{\small#3} & \textit{\small #4} \\
    \end{tabular*}\vspace{-6pt}
}

\newcommand{\resumeSubItem}[2]{\resumeItem{#1}{#2}\vspace{-5pt}}

\renewcommand{\labelitemii}{$\circ$}

\newcommand{\resumeSubHeadingListStart}{\begin{itemize}[leftmargin=*]}
\newcommand{\resumeSubHeadingListEnd}{\end{itemize}\vspace{-8pt}}
\newcommand{\resumeItemListStart}{\begin{itemize}}
\newcommand{\resumeItemListEnd}{\end{itemize}\vspace{-6pt}}


%-------------------------------------------
%%%%%%  CV STARTS HERE  %%%%%%%%%%%%%%%%%%%%%%%%%%%%


\begin{document}

%----------HEADING-----------------
\begin{tabular*}{\textwidth}{l@{\extracolsep{\fill}}r}
  \textbf{\Large Sırrı Kaan Oğuzkan} & E-posta : \href{mailto:kaan.oguzkan@ug.bilkent.edu.tr}{kaan.oguzkan@ug.bilkent.edu.tr}\\
  \href{https://www.linkedin.com/in/kaan-oguzkan/}{{linkedin.com/in/kaan-oguzkan}} & \href{https://kaanoguzkan.com}{{kaanoguzkan.com}} \\
  \href{https://github.com/kaanoguzkan}{{github.com/kaanoguzkan}} & \href{https://orcid.org/0009-0000-3272-7333}{{ORCID: 0009-0000-3272-7333}} \\
\end{tabular*}

\section{Eğitim}
  \resumeSubHeadingListStart
    \resumeSubheading
      {Bilkent Üniversitesi}{Ankara, Türkiye}
      {Bilgisayar Mühendisliği Lisans}{Eylül 2023 -- Devam ediyor}
  \resumeSubHeadingListEnd
%-----------EXPERIENCE-----------------
\section{Deneyim}
  \resumeSubHeadingListStart

    \resumeSubheading
  {JotForm}{Ankara, Türkiye}
  {Yazılım Mühendisliği Stajyeri}{Haziran 2026 - Ağustos 2026}
  \resumeItemListStart
      \resumeItem{Element Registry Mimarisi}
        {Sürükle-bırak bir e-posta oluşturucu (\textbf{React 18}, \textbf{TypeScript}, \textbf{Redux}) için eklenti-sözleşmesi (plugin-contract) mimarisi tasarladım; \textbf{20} içerik bloğu türü canvas, panel, önizleme ve dışa aktarıcı boyunca tek bir tipli arayüzü paylaşıyor, böylece yeni bloklar mevcut tüketicilerin \textbf{hiçbirine} dokunmayı gerektirmedi.}
      \resumeItem{Toplu Gönderim Hattı}
        {Asenkron bir toplu e-posta hattı (PHP, MySQL, Redis) geliştirdim; parçalı SMTP oturumlarıyla (\textbf{25} alıcı/grup) teslimat yapan bir cron worker ve takılı kalan gruplar için \textbf{900 sn} zaman aşımına dayalı geri alma (reclaim) mekanizması ekledim.}
    \resumeItemListEnd

    \resumeSubheading
  {Alkan Biyoenformatik ve Hesaplamalı Genomik Laboratuvarı}{Ankara, Türkiye}
  {Lisans Gönüllü Araştırmacı}{Şubat 2026 - Devam ediyor}
  \resumeItemListStart
      \resumeItem{CuCLARK GPU \& CUDA Modernizasyonu}
        {CuCLARK'ı (GPU tabanlı metagenomik sınıflandırıcı, orijinali CUDA 7.5 / tek mimari) forkladım ve modernize ettim; \textbf{CUDA 11.8}'e yükselttim ve \textbf{10} GPU neslini (Kepler \texttt{sm\_35}'ten Hopper \texttt{sm\_90}'a; K80, V100, A100, RTX 30xx/40xx, H100) kapsayan çoklu mimari desteği ekledim.}
      \resumeItem{Build Sistemi \& Konteynerleştirme}
        {Eski \texttt{install.sh} betiğini \textbf{CMake 3.20} tabanlı bir build (C++17, CUDA 14, isteğe bağlı OpenMP) ile değiştirdim ve \texttt{alkanlab/cuclark} olarak yayımlanan çok mimarili bir \texttt{Dockerfile} yazdım; orijinal betikleri saran \textbf{6} alt komutlu birleşik bir CLI ekledim.}
      \resumeItem{Toplu Sınıflandırma İşlem Hattı}
        {Küçük girdilerde tekrarlanan okuma (read) atamalarına yol açan bir bayt-ofset batch hatasını teşhis ettim; devam ettirme, çökme sonrası kara liste ve ilerleme takibi içeren bir liste modu işlem hattı kurarak \textbf{3000+} FASTA dosyasını tek bir veritabanı yüklemesinde sınıflandırdım.}
    \resumeItemListEnd

    \resumeSubheading
  {Look \& Cash}{Ankara, Türkiye}
  {Yarı Zamanlı Yazılım Mühendisi}{Şubat 2025 - Ekim 2025}
  \resumeItemListStart
      \resumeItem{Bilet Tabanlı Destek Sistemi}
        {Web ve mobil için bilet (ticket) iş akışı geliştirdim (öncelik, yorumlar, ekler, durumlar). \textbf{13} REST uç noktası, \textbf{4} bilet durumu ve \textbf{5} yönetici aksiyonu sevk ettim; toplu işlemler ve gelişmiş filtreleme ile ortalama yönetici işlem süresini \textbf{\%25} azalttım.}
      \resumeItem{Mobil Gerçek Zamanlı Destek Sohbeti}
        {Mobilde, backend kalıcılığı ve durum senkronizasyonu ile gerçek zamanlı bilet sohbeti geliştirdim; \textbf{150} eşzamanlı oturum altında \textbf{300 ms} p95 iletim süresine ve \textbf{\%99,5} başarıya ulaştım.}
      \resumeItem{Raporlama \& Grafik Üretimi}
        {MongoDB üzerinde tarih aralığı ve kampanya filtreleri ile KPI raporlama API'leri (biletler, kampanyalar, taramalar, bildirimler) geliştirdim. \textbf{4} grafik tipini besleyen \textbf{14} pano veri seti ürettim; rapor üretim süresini indeksleme ile \textbf{6,2 sn}'den \textbf{1,9 sn}'ye düşürdüm.}
    \resumeItemListEnd

  \resumeSubHeadingListEnd


%-----------PROJECTS-----------------
\section{Projeler}
  \resumeSubHeadingListStart

\resumeSubItem{osguide.org — Etkileşimli İşletim Sistemi Öğrenme Platformu}
{Bilkent'in CS342 müfredatı etrafında yapılandırılmış etkileşimli bir platform geliştirdim (\textbf{React}, \textbf{TypeScript}, \textbf{Vite}).}
\resumeItemListStart
  \resumeItem{Kavramdan Simülatöre Katmanlar}
    {Her konu için üç katmanlı bir öğrenme akışı tasarladım — kavram anlatımı, animasyon ve etkileşimli simülatör — zamanlama, bellek yönetimi ve senkronizasyon gibi temel işletim sistemi konularını kapsıyor.}
\resumeItemListEnd

\resumeSubItem{NutriApp — Yemek Öneri Platformu}
{Amazon University Engagement Program kapsamında bir yemek planlama uygulamasının backend servislerini geliştirdim.}
\resumeItemListStart
  \resumeItem{Sunucusuz REST API'leri}
    {AWS Lambda üzerinde \textbf{16} REST uç noktası geliştirdim; \textbf{People's Choice Award} kazandım, \textbf{7 takım arasında 1.} oldum.}
\resumeItemListEnd

  \resumeSubHeadingListEnd


%-----------SKILLS-----------------
\section{Yetkinlikler}
 \resumeSubHeadingListStart
    \item{
      \textbf{Programlama Dilleri}{: C++, TypeScript, JavaScript, PHP, Java}
    }
    \item{
      \textbf{Çerçeveler \& Araçlar}{: React, Redux, Node.js, AWS Lambda, CUDA, Docker, Git}
    }
    \item{
      \textbf{Veritabanları}{: MongoDB, DynamoDB, MySQL, Redis}
    }

 \resumeSubHeadingListEnd


\end{document}
